\subsubsection{Mô hình Probit}

Thay vì sử dụng hàm sigmoid, mô hình Probit định nghĩa xác suất thuộc lớp dương
thông qua hàm phân phối tích lũy (CDF) của phân phối Gaussian chuẩn:
\begin{equation}
  p(C_1 \mid \boldsymbol{x}) = \Phi(\boldsymbol{w}^{\intercal}\boldsymbol{x} +
  w_0),
\end{equation} trong đó hàm \(\Phi(\cdot)\) được định nghĩa bởi:
\begin{equation}
  \Phi(a) = \int_{-\infty}^{a} \mathcal{N}(t \mid 0,1)\, dt.
\end{equation}

Mô hình Probit có thể được diễn giải từ mô hình biến ẩn (latent variable
model). Cụ thể, ta giả sử tồn tại một biến ẩn \(z\) được sinh bởi:
\begin{equation}
  z = \boldsymbol{w}^{\intercal}\boldsymbol{x} + w_0 + \epsilon,
\end{equation} trong đó nhiễu \(\epsilon \sim \mathcal{N}(0,1)\). Khi đó, nhãn
quan sát được xác định theo quy tắc: \begin{equation}
  t = \begin{cases}
    1 & \text{nếu $z > 0$}, \\
    0 & \text{ngược lại}.
  \end{cases}
\end{equation} Từ đó suy ra xác suất có dạng: \begin{equation}
  p(t = 1 \mid \boldsymbol{x}) = \mathbb{P}(z > 0)
  = \Phi(\boldsymbol{w}^{\intercal}\boldsymbol{x} + w_0).
\end{equation}

Nhờ cấu trúc dựa trên Gaussian noise, mô hình Probit thường phù hợp hơn hồi quy
Logistic trong các trường hợp dữ liệu có nhãn bị nhiễu hoặc có giả định sai số
tuân theo phân phối chuẩn.

\paragraph{So sánh với hồi quy Logistic}

Phần này sẽ so sánh thực nghiệm của mô hình Probit với hồi quy Logistic
(sigmoid). Dữ liệu được sử dụng là dữ liệu được sinh bằng hàm
\texttt{make\_blobs} với 1000 điểm dữ liệu, tâm là $(-2, 0)$ và $(0, 1.5)$. Dữ
liệu sau đó được biến đổi tuyến tính qua ma trận: \begin{equation}
  A = \begin{bmatrix}
    0.4 & 0.2\\
    -0.2 & 0.5
  \end{bmatrix}.
\end{equation} Dữ liệu được biểu diễn ở \figurename~\ref{fig:probit-data}.

\subparagraph{Đường biên quyết định}

\figurename~\ref{fig:sgd-decision-boundary},
\ref{fig:sigmoid-decision-boundary}, và~\ref{fig:probit-decision-boundary}
tương ứng là đường biên quyết định tương ứng của ba mô hình hồi quy Logistic
(scikit-learn), hồi quy Logistic (tự cài đặt) và Probit.

\begin{figure}[htbp]
  \centering
  \includegraphics[width=0.9\textwidth]{probit-data}
  \caption{Dữ liệu tạo sinh cho phần mô hình Probit và xấp xỉ Laplace}
  \label{fig:probit-data}
\end{figure}

\begin{figure}[htbp]
  \centering
  \includegraphics[width=0.9\textwidth]{sgd-decision-boundary}
  \caption{Đường biên quyết định của mô hình SGDClassifier (scikit-learn)}
  \label{fig:sgd-decision-boundary}
\end{figure}

\begin{figure}[htbp]
  \centering
  \includegraphics[width=0.9\textwidth]{sigmoid-decision-boundary}
  \caption{Đường biên quyết định của mô hình hồi quy Logistic}
  \label{fig:sigmoid-decision-boundary}
\end{figure}

\begin{figure}[htbp]
  \centering
  \includegraphics[width=0.9\textwidth]{probit-decision-boundary}
  \caption{Đường biên quyết định của mô hình Probit}
  \label{fig:probit-decision-boundary}
\end{figure}

Ta thấy rằng đường biên của 3 mô hình sau khi hội tụ là khá giống nhau. Trong
đó, mô hình Probit có đường biên có vẻ dốc hơn.

\subparagraph{Xác suất đầu ra}

\figurename~\ref{fig:sigmoid-probability-output}
và~\ref{fig:probit-probability-output} tương ứng là xác suất đầu ra của hồi quy
Logistic và Probit.

\begin{figure}[htbp]
  \centering
  \includegraphics[width=0.9\textwidth]{sigmoid-probability-output}
  \caption{Xác suất đầu ra của mô hình hồi quy Logistic}
  \label{fig:sigmoid-probability-output}
\end{figure}

\begin{figure}[htbp]
  \centering
  \includegraphics[width=0.9\textwidth]{probit-probability-output}
  \caption{Xác suất đầu ra của mô hình Probit}
  \label{fig:probit-probability-output}
\end{figure}

Quan sát các biểu đồ xác suất đầu ra của mô hình hồi quy Logistic và Probit
cho thấy sự khác biệt chủ yếu nằm ở vùng chuyển tiếp xác suất quanh đường biên
quyết định, tức khu vực mà xác suất thay đổi từ gần 0 sang gần 1.

Trong cả hai mô hình, vùng này tương ứng với miền giá trị của $z =
\boldsymbol{w}^\intercal \boldsymbol{x}$ lân cận 0, nơi hàm liên kết có độ dốc
lớn nhất. Tuy nhiên, hình dạng của vùng chuyển tiếp có sự khác biệt nhẹ: mô
hình Sigmoid có xu hướng phân bố sự thay đổi xác suất trên một khoảng giá trị
$z$ rộng hơn, trong khi mô hình Probit thể hiện sự tập trung biến thiên xác
suất mạnh hơn quanh vùng trung tâm.

Điều này có thể được giải thích từ dạng đạo hàm của hai hàm liên kết. Cụ thể,
Sigmoid có gradient phụ thuộc vào chính giá trị đầu ra, trong khi Probit có
gradient tương ứng với mật độ của phân phối Gaussian chuẩn. Do đó, Probit có xu
hướng tạo ra sự thay đổi xác suất rõ rệt hơn gần $z = 0$, nhưng giảm nhanh
hơn khi đi xa khỏi vùng trung tâm.

Tuy nhiên, cần lưu ý rằng sự khác biệt về vùng chuyển tiếp giữa hai mô hình
không hoàn toàn cố định, mà phụ thuộc vào tham số học được trong quá trình tối
ưu. Trên thực nghiệm, cả hai mô hình thường cho đường biên quyết định tương đối
tương đồng, và khác biệt chủ yếu thể hiện ở cách phân bố xác suất cục bộ quanh
decision boundary.

\subparagraph{Độ nhạy với nhãn nhiễu}

\tablename~\ref{tab:probit-sigmoid-label-noise-sensitivity-part1}
và~\ref{tab:probit-sigmoid-label-noise-sensitivity-part2} cho ta thấy độ nhạy
với nhãn nhiễu của hai mô hình Probit và hồi quy Logistic. Tỉ lệ dữ liệu bị
thay đổi nhãn là $0.2$. Trong đó xác suất bị thay đổi tỉ lệ với khoảng cách tới
đường biên quyết định của model baseline (\texttt{SGDClassifier} của
scikit-learn).

\begin{table}[htbp]
  \centering
  \caption{Độ nhạy với nhãn nhiễu trên 2 mô hình Probit và hồi quy tuyến tính
  (phần 1)}
  \label{tab:probit-sigmoid-label-noise-sensitivity-part1}
  \import{\tablespath}{probit-sigmoid-label-noise-sensitivity-part1}
\end{table}

\begin{table}[htbp]
  \centering
  \caption{Độ nhạy với nhãn nhiễu trên 2 mô hình Probit và hồi quy tuyến tính
  (phần 2)}
  \label{tab:probit-sigmoid-label-noise-sensitivity-part2}
  \import{\tablespath}{probit-sigmoid-label-noise-sensitivity-part2}
\end{table}

Ta thấy rằng Probit hầu như không thay đổi hiệu suất trên các chỉ số. Điều này
cho thấy mô hình này ổn định hơn với nhãn nhiễu. Tuy nhiên sự sai khác so với
hồi quy tuyến tính khá nhỏ nên cũng chưa có ý nghĩa rõ ràng.
